% Copyright (c) 2019 FLAG lab
% Implementation chapter

\chapter{Implementation}\label{chap:implementation}

\section{Introduction}

This chapter describes the concrete implementation of the CRDT framework in Athena, showing how the theoretical foundations from Chapter~\ref{chap:theory} are realized as executable code.

\section{Module Organization}

The implementation is organized into several key modules:

\subsection{Framework Layer}

The framework provides the core infrastructure:

\begin{itemize}
  \item \texttt{payload.ath} -- Unified payload framework with \texttt{PayloadOps} interface and example implementations (TextPayload, GCounterPayload, GSetPayload)
  \item \texttt{backbone.ath} -- Unified backbone with distributed system semantics, axiomatic reduce operations, and vector clock integration
  \item \texttt{vector\_clocks.ath} -- Unified vector clock module with formal proofs of semilattice properties (idempotence, commutativity, associativity)
  \item \texttt{nat-max.ath} -- Natural number maximum with complete formal proofs
\end{itemize}

\subsection{Theory Layer}

Formal foundations for verification:

\begin{itemize}
  \item \texttt{theory/semilattice\_theory\_v2.ath} -- Abstract semilattice with 5 axioms and 3 proven theorems (absorption, monotone-left, monotone-right)
  \item \texttt{theory/gcounter\_theory\_v3.ath} -- Canonical example showing complete adapter pattern with theorem inheritance
\end{itemize}

\subsection{CRDT Implementations}

Concrete CRDT types using the semilattice adapter pattern:

\begin{itemize}
  \item \textbf{GCounter} (\texttt{crdts/gcounter\_semilattice\_adapter.ath}) -- Grow-only counter using \texttt{N.Max}
  \item \textbf{GSet} (\texttt{crdts/gset\_semilattice\_adapter.ath}) -- Grow-only set using set union
  \item \textbf{PN-Counter} (\texttt{crdts/pn\_counter\_semilattice\_adapter.ath}) -- Product of two GCounters
  \item \textbf{2P-Set} (\texttt{crdts/2p\_set\_semilattice\_adapter.ath}) -- Product of two GSets
  \item \textbf{OR-Set} (\texttt{crdts/orset\_semilattice\_adapter.ath}) -- Observed-remove set with causal consistency
  \item \textbf{MV-Register} (\texttt{crdts/mv\_register\_semilattice\_adapter.ath}) -- Multi-value register keeping concurrent writes
  \item \textbf{LWW-Register} (\texttt{crdts/lww\_register\_semilattice\_adapter.ath}) -- Last-writer-wins with timestamps
\end{itemize}

\section{Payload Framework}

The payload framework (\texttt{src/framework/payload.ath}) provides a generic interface for CRDT operations through the \texttt{PayloadOps} structure. Each CRDT implements three core operations:

\begin{itemize}
  \item \texttt{payload-init} -- Returns the initial (bottom) state
  \item \texttt{payload-step} -- Applies an update operation to the current state
  \item \texttt{payload-join} -- Merges two replica states
\end{itemize}

These operations correspond to the abstract CRDT specification from Chapter~\ref{chap:theory}:

\begin{center}
\begin{tabular}{l|l}
\textbf{Abstract Theory} & \textbf{Payload Implementation} \\
\hline
\texttt{initial} & \texttt{payload-init ops} \\
\texttt{update(s, v)} & \texttt{payload-step ops v s} \\
\texttt{merge(s1, s2)} & \texttt{payload-join ops s1 s2} \\
\end{tabular}
\end{center}

\subsection{GCounter Implementation}

The GCounter is implemented using the semilattice adapter pattern in \texttt{gcounter\_semilattice\_adapter.ath}. It uses natural numbers (\texttt{N}) as the state sort with \texttt{N.Max} as the join operation.

\AthenaCode{lst:gcounter-impl}{GCounter semilattice adapter}
\begin{lstlisting}[language=Athena,style=athena]
module GCounterAdapter {
  open N
  
  # Define concrete operations
  define (gc-join x y) := (N.Max x y)
  define gc-initial := zero
  
  # Prove 5 semilattice axioms
  conclude gc-commutative :=
    (forall ?n1:N ?n2:N . (gc-join ?n1 ?n2) = (gc-join ?n2 ?n1))
    pick-any n1:N n2:N
      (!chain [(gc-join n1 n2)
             = (N.Max n1 n2)       [definition]
             = (N.Max n2 n1)       [N.Max.commutative]
             = (gc-join n2 n1)     [definition]])
  
  # ... (associative, idempotent, left/right identity)
  
  # Create renaming adapter
  define GCounter-Adapter :=
    (renaming |{SemilatticeTheory.join := gc-join,
                SemilatticeTheory.bot := gc-initial}|)
  
  # Inherit convergence theorem automatically
  define gc-convergence :=
    (get-property SemilatticeTheory.absorption
                  GCounter-Adapter
                  SL-Theory)
}
\end{lstlisting}

The implementation uses Athena's \textbf{theory/adapter pattern} (Chapter 14): define concrete operations, prove the five semilattice axioms using \texttt{!chain} proofs, create a renaming adapter, and automatically inherit convergence theorems. This pattern is used for all 7 verified CRDTs.

\section{Backbone Integration}

The \texttt{backbone.ath} module integrates payload operations with vector clocks to ensure causal consistency through the \texttt{VCBackbone} integration:

\begin{itemize}
  \item \textbf{Configuration}: Contains actor states (\texttt{FMap Ide State}), vector clock registry (\texttt{Registry}), and pending messages (\texttt{MSet Message})
  \item \textbf{Clock operations}: \texttt{cfg-tick} increments local clock, \texttt{clock-merge-entry} updates registry on message delivery
  \item \textbf{Message delivery}: \texttt{reduce} processes one message, updating both state (via \texttt{payload-apply}) and clock (via \texttt{clock-join})
  \item \textbf{Termination}: \texttt{reduce*} drains all messages until network is empty
\end{itemize}

The vector clock module provides \textbf{proven semilattice properties}:

\begin{itemize}
  \item \texttt{clock-join-lookup-is-max} -- Proven by induction that clock join is pointwise max
  \item \texttt{idempotent-pointwise}, \texttt{commutative-pointwise}, \texttt{associative-pointwise} -- Proven from \texttt{N.Max} properties
  \item \texttt{ClockJoin.idempotent-proven} -- Lifted to full clock equality via structural induction
\end{itemize}

This integration ensures that when updates occur:

\begin{enumerate}
  \item The local vector clock is incremented (\texttt{cfg-tick})
  \item The payload state is updated via \texttt{payload-apply}
  \item A message containing the operation and timestamp is broadcast
  \item Recipients deliver messages, merging both state and clocks
  \item Convergence is guaranteed by semilattice properties of both payload and clocks
\end{enumerate}

\section{Executable Testing}

All CRDT implementations are executable and can be tested directly in Athena:

\begin{lstlisting}[language=Athena,style=athena]
# Create initial GCounter state
define s0 := (payload-initial-state GCOUNTER_PAYLOAD)

# Apply increments from different actors
define s1 := (payload-step GCOUNTER_PAYLOAD (inc 'a one) s0)
define s2 := (payload-step GCOUNTER_PAYLOAD (inc 'b two) s1)

# Query the total value
> (gcounter-total s2)
three
\end{lstlisting}

This executable nature allows for rapid prototyping and testing of CRDT properties.

\section{Relationship to Theory}

The implementations use the \textbf{semilattice adapter pattern} (Athena Chapter 14) to formally connect executable code to mathematical theory:

\subsection{Pattern Structure}

For each CRDT, we follow four steps:

\begin{enumerate}
  \item \textbf{Define concrete operations}: Map abstract join/bot to concrete operations (e.g., \texttt{gc-join := N.Max})
  \item \textbf{Prove axioms}: Use \texttt{!chain} proofs to show concrete operations satisfy the 5 semilattice axioms
  \item \textbf{Create adapter}: Use \texttt{renaming} to bind abstract theory to concrete implementation
  \item \textbf{Inherit theorems}: Automatically get convergence theorems (absorption, monotonicity) via \texttt{get-property}
\end{enumerate}

\subsection{Verified Properties}

Each CRDT has \textbf{machine-checked proofs} of:

\begin{itemize}
  \item \texttt{commutative} -- \texttt{join(s1, s2) = join(s2, s1)}
  \item \texttt{associative} -- \texttt{join(join(s1, s2), s3) = join(s1, join(s2, s3))}
  \item \texttt{idempotent} -- \texttt{join(s, s) = s}
  \item \texttt{left-identity} -- \texttt{join(bot, s) = s}
  \item \texttt{right-identity} -- \texttt{join(s, bot) = s}
\end{itemize}

These five properties \textbf{guarantee Strong Eventual Consistency}: any two replicas that have received the same set of updates (in any order, with any duplicates) will converge to the same state via repeated join operations.

\subsection{Product Semilattices}

Complex CRDTs (PN-Counter, 2P-Set) use the \textbf{product semilattice} principle: if $(S_1, \sqcup_1, \bot_1)$ and $(S_2, \sqcup_2, \bot_2)$ are semilattices, then $(S_1 \times S_2, \text{pointwise-}\sqcup, (\bot_1, \bot_2))$ is also a semilattice. This enables compositional verification:

\begin{itemize}
  \item PN-Counter = GCounter $\times$ GCounter (verify GCounter once, get PN-Counter for free)
  \item 2P-Set = GSet $\times$ GSet (verify GSet once, get 2P-Set for free)
\end{itemize}

\subsection{Executable and Verified}

The adapter pattern provides both:

\begin{itemize}
  \item \textbf{Formal correctness} -- \texttt{!chain} proofs guarantee semilattice properties → SEC
  \item \textbf{Practical usability} -- Executable operations can be tested in demos
\end{itemize}

All 7 CRDTs (GCounter, GSet, PN-Counter, 2P-Set, OR-Set, MV-Register, LWW-Register) follow this pattern, totaling 1,603 lines of verified code with machine-checked convergence guarantees.
